\documentclass{homework}
\author{Mufan Shi 23307130075}
\class{Introduction to Robotics}
\date{\today}
\title{Homework of Chapter3}

\usepackage{ctex} 
\usepackage{subcaption}
\graphicspath{{./media/}}

\begin{document} \maketitle

\question 

\begin{sol}

\img<q1>[0.5]{Question 1}{figs/q1.png}

\tbl<tb1>{DH-table for question 1} {
    i & $ a_{i-1} $ & $ \alpha_{i-1} $ & $ d_i $ & $ \theta_i $ \\
    1 & 0 & 0 & $ L_1 $ & $ \theta_1 $ \\
    2 & 0 & $ -\frac{\pi}{2} $ & 0 & $ \theta_2 $ \\
    3 & $ L_2 $ & $ -\frac{\pi}{2} $ & $ d_3 $ & 0 \\
  }
  
\end{sol}

\question 

\begin{sol}

\img<q2>[0.5]{Question 2}{figs/q2.png}

\tbl<tb2>{DH-table for question 2} {
    i & $ a_{i-1} $ & $ \alpha_{i-1} $ & $ d_i $ & $ \theta_i $ \\
    1 & 0 & 0 & $ L_1 $ & $ \theta_1 $ \\
    2 & 0 & $ -\frac{\pi}{2} $ & 0 & $ \theta_2 $ \\
    3 & $ L_2 $ & 0 & 0 & $ \theta_3 $ \\
  }
  
\end{sol}

\question 

\begin{sol}

我们把末端执行器作为第5个关节，并把它当作一个固定关节来处理。

\img<q3>[0.75]{Question 3}{figs/q3.png}

\tbl<tb3>{DH-table for question 3} {
    i & $ a_{i-1} $ & $ \alpha_{i-1} $ & $ d_i $ & $ \theta_i $ \\
    1 & 0 & 0 & 0 & $ \theta_1 $ \\
    2 & $ l_1 $ & 0 & 0 & $ \theta_2 $ \\
    3 & $ l_2 $ & 0 & 0 & $ \theta_3 $ \\
    4 & $ l_3 $ & 0 & 0 & $ \theta_4 $ \\
    5 & $ l_4 $ & 0 & 0 & 0 \\
  }

求出各坐标系变换矩阵如下：

\[
  _1^0\mathbf{T} =
  \begin{bmatrix}
    \cos\theta_1 & -\sin\theta_1 & 0 & 0 \\
    \sin\theta_1 & \cos\theta_1 & 0 & 0 \\
    0 & 0 & 1 & 0 \\
    0 & 0 & 0 & 1
  \end{bmatrix}
\]

\[
  _2^1\mathbf{T} =
  \begin{bmatrix}
    \cos\theta_2 & -\sin\theta_2 & 0 & l_1 \\
    \sin\theta_2 & \cos\theta_2 & 0 & 0 \\
    0 & 0 & 1 & 0 \\
    0 & 0 & 0 & 1
  \end{bmatrix}
\]

\[
  _3^2\mathbf{T} =
  \begin{bmatrix}
    \cos\theta_3 & -\sin\theta_3 & 0 & l_2 \\
    \sin\theta_3 & \cos\theta_3 & 0 & 0 \\
    0 & 0 & 1 & 0 \\
    0 & 0 & 0 & 1
  \end{bmatrix}
\]

\[
  _4^3\mathbf{T} =
  \begin{bmatrix}
    \cos\theta_4 & -\sin\theta_4 & 0 & l_3 \\
    \sin\theta_4 & \cos\theta_4 & 0 & 0 \\
    0 & 0 & 1 & 0 \\
    0 & 0 & 0 & 1
  \end{bmatrix}
\]

\[
  _5^4\mathbf{T} =
  \begin{bmatrix}
    1 & 0 & 0 & l_4 \\
    0 & 1 & 0 & 0 \\
    0 & 0 & 1 & 0 \\
    0 & 0 & 0 & 1
  \end{bmatrix}
\]

末端执行器在基坐标下的位姿为：

\[
  _5^0\mathbf{T} = _1^0\mathbf{T} _2^1\mathbf{T} _3^2\mathbf{T} _4^3\mathbf{T} _5^4\mathbf{T} =
  \begin{bmatrix}
    c_{1234} & -s_{1234} & 0 & l_1 c_1 + l_2 c_{12} + l_3 c_{123} + l_4 c_{1234} \\
    s_{1234} & c_{1234} & 0 & l_1 s_1 + l_2 s_{12} + l_3 s_{123} + l_4 s_{1234} \\
    0 & 0 & 1 & 0 \\
    0 & 0 & 0 & 1
  \end{bmatrix}
\]
  
\end{sol}

\question 

\begin{sol}

修正后的DH表如下：
\tbl<tb4>{DH-table for question 4} {
    i & $ a_{i-1} $ & $ \alpha_{i-1} $ & $ d_i $ & $ \theta_i $ \\
    1 & 0 & 0 & 0 & $ \theta_1 $ \\
    2 & 0 & $ \frac{\pi}{2} $ & $ b_1 + b_2 $ & $ \theta_2 $ \\
    3 & 8 & 0 & 0 & $ \theta_3 $ \\
    4 & 0 & $ \frac{\pi}{2} $ & 8 & $ \theta_4 $ \\
    5 & 0 & $ -\frac{\pi}{2} $ & 0 & $ \theta_5 $ \\
    6 & 0 & $ -\frac{\pi}{2} $ & $ b_3 $ & $ \theta_6 $ 
  }

各连杆变换矩阵如下：

\[
  _1^0\mathbf{T} =
  \begin{bmatrix}
    \cos\theta_1 & -\sin\theta_1 & 0 & 0 \\
    \sin\theta_1 & \cos\theta_1 & 0 & 0 \\
    0 & 0 & 1 & 0 \\
    0 & 0 & 0 & 1
  \end{bmatrix}
\]

\[
  _2^1\mathbf{T} =
  \begin{bmatrix}
    \cos\theta_2 & -\sin\theta_2 & 0 & 0 \\
    0 & 0 & -1 & - b_1 - b_2 \\
    \sin\theta_2 & \cos\theta_2 & 0 & 0 \\
    0 & 0 & 0 & 1
  \end{bmatrix}
\]

\[
  _3^2\mathbf{T} =
  \begin{bmatrix}
    \cos\theta_3 & -\sin\theta_3 & 0 & 8 \\
    \sin\theta_3 & \cos\theta_3 & 0 & 0 \\
    0 & 0 & 1 & 0 \\
    0 & 0 & 0 & 1
  \end{bmatrix}
\]

\[
  _4^3\mathbf{T} =
  \begin{bmatrix}
    \cos\theta_4 & -\sin\theta_4 & 0 & 0 \\
    0 & 0 & -1 & -8 \\
    \sin\theta_4 & \cos\theta_4 & 0 & 0 \\
    0 & 0 & 0 & 1
  \end{bmatrix}
\]

\[
  _5^4\mathbf{T} =
  \begin{bmatrix}
    \cos\theta_5 & -\sin\theta_5 & 0 & 0 \\
    0 & 0 & 1 & 0 \\
    -\sin\theta_5 & -\cos\theta_5 & 0 & 0 \\
    0 & 0 & 0 & 1
  \end{bmatrix}
\]

\[
  _6^5\mathbf{T} =
  \begin{bmatrix}
    \cos\theta_6 & -\sin\theta_6 & 0 & 0 \\
    0 & 0 & 1 & b_3 \\
    -\sin\theta_6 & -\cos\theta_6 & 0 & 0 \\
    0 & 0 & 0 & 1
  \end{bmatrix}
\]

从DH表我们可以看出，关节2和关节3是平行的，所以 $ _2^1\mathbf{T} $ 和 $ _3^2\mathbf{T} $ 的乘积用和差化积公式将得到一个简化的表达式。
\[
  _3^1\mathbf{T} = _2^1\mathbf{T} _3^2\mathbf{T} =
  \begin{bmatrix}
    c_{23} & -s_{23} & 0 & 8c_2 \\
    0 & 0 & -1 & - b_1 - b_2 \\
    s_{23} & c_{23} & 0 & 8s_2 \\
    0 & 0 & 0 & 1
  \end{bmatrix}
\]

末端执行器的变换矩阵为：

\[
  _6^0\mathbf{T} = _1^0\mathbf{T} _3^1\mathbf{T} _4^3\mathbf{T} _5^4\mathbf{T} _6^5\mathbf{T} =
  \begin{bmatrix} 
r_{11} & r_{12} & r_{13} & p_x \\ 
r_{21} & r_{22} & r_{23} & p_y \\ 
r_{31} & r_{32} & r_{33} & p_z \\ 
0 & 0 & 0 & 1
\end{bmatrix}
\] 

\begin{align*} 
r_{11} &= c_{6}\,\left(c_{5}\,\left(s_{1}\,s_{4}+c_{23}\,c_{1}\,c_{4}\right)-s_{23}\,c_{1}\,s_{5}\right)-s_{6}\,\left(c_{4}\,s_{1}-c_{23}\,c_{1}\,s_{4}\right) \\ 
r_{12} &= -s_{6}\,\left(c_{5}\,\left(s_{1}\,s_{4}+c_{23}\,c_{1}\,c_{4}\right)-s_{23}\,c_{1}\,s_{5}\right)-c_{6}\,\left(c_{4}\,s_{1}-c_{23}\,c_{1}\,s_{4}\right) \\ 
r_{13} &= -s_{5}\,\left(s_{1}\,s_{4}+c_{23}\,c_{1}\,c_{4}\right)-s_{23}\,c_{1}\,c_{5} \\ 
p_x &= 8\,c_{1}\,c_{2}+s_{1}\,\left(b_{1}+b_{2}\right)-b_{3}\,\left(s_{5}\,\left(s_{1}\,s_{4}+c_{23}\,c_{1}\,c_{4}\right)+s_{23}\,c_{1}\,c_{5}\right)+8\,s_{23}\,c_{1} \\ 
r_{21} &= s_{6}\,\left(c_{1}\,c_{4}+c_{23}\,s_{1}\,s_{4}\right)-c_{6}\,\left(c_{5}\,\left(c_{1}\,s_{4}-c_{23}\,c_{4}\,s_{1}\right)+s_{23}\,s_{1}\,s_{5}\right) \\ 
r_{22} &= s_{6}\,\left(c_{5}\,\left(c_{1}\,s_{4}-c_{23}\,c_{4}\,s_{1}\right)+s_{23}\,s_{1}\,s_{5}\right)+c_{6}\,\left(c_{1}\,c_{4}+c_{23}\,s_{1}\,s_{4}\right) \\ 
r_{23} &= s_{5}\,\left(c_{1}\,s_{4}-c_{23}\,c_{4}\,s_{1}\right)-s_{23}\,c_{5}\,s_{1} \\ 
p_y &= 8\,s_{23}\,s_{1}+8\,c_{2}\,s_{1}-c_{1}\,\left(b_{1}+b_{2}\right)+b_{3}\,\left(s_{5}\,\left(c_{1}\,s_{4}-c_{23}\,c_{4}\,s_{1}\right)-s_{23}\,c_{5}\,s_{1}\right) \\ 
r_{31} &= c_{6}\,\left(c_{23}\,s_{5}+s_{23}\,c_{4}\,c_{5}\right)+s_{23}\,s_{4}\,s_{6} \\ 
r_{32} &= s_{23}\,c_{6}\,s_{4}-s_{6}\,\left(c_{23}\,s_{5}+s_{23}\,c_{4}\,c_{5}\right) \\ 
r_{33} &= c_{23}\,c_{5}-s_{23}\,c_{4}\,s_{5} \\ 
p_z &= 8\,s_{2}-8\,c_{23}+b_{3}\,\left(c_{23}\,c_{5}-s_{23}\,c_{4}\,s_{5}\right)
\end{align*}
  
\end{sol}

\question 

\begin{sol}

公垂线的方向向量应该同时垂直于 $ \hat{m} $ 和 $ \hat{n} $，所以我们可以用叉积得到公垂线的方向向量：$ \hat{d} = \hat{m} \times \hat{n} $。公垂线的长度 a 是两条直线之间的最短距离。这个最短距离可以通过向量 $ \mathbf{p} - \mathbf{q} $ 在公垂线方向上的投影的绝对值得到，也就是

\[
  a = \frac{(\mathbf{p} - \mathbf{q})(\hat{m} \times \hat{n})}{|\hat{m} \times \hat{n}|}
\]

当两条直线平行时，$ \hat{m} \times \hat{n} = \mathbf{0} $，上述表达式分母为零。但此时公垂线长度 a 为点 q 到直线 $ \mathbf{p} + t \hat{m} $ 的距离，即：

\[
  a = |(\mathbf{p} - \mathbf{q}) \times \hat{m}|
\]

需要注意的是，当两条直线相交时，公垂线长度 $ a = 0 $。

角 $ \alpha $ 是两条直线在与公垂线相垂直的平面上的投影的夹角。由于公垂线垂直于两条直线，两条直线的方向矢量在垂直于公垂线的平面上的投影即为自身，因此 $ \alpha $ 是方向矢量 $ \hat{m} $ 和 $ \hat{n} $ 之间的最小夹角。表达式为：

\[
  \alpha = \arccos (|\hat{m} \cdot \hat{n}|)
\]
  
\end{sol}

\question 

\begin{sol}

这个问题其实取决于坐标系的放置方式。如果坐标系按照 DH 约定放置（即 $ x_{i-1} $ 轴沿 $ z_{i-1} $ 和 $ z_i $ 的公垂线方向），则任意刚体变换都可以用这四个参数表示。这种放置方式确保了变换矩阵具有上述形式，从而用四个参数捕获了变换的六个自由度（三个平移和三个旋转），因为坐标系的约束减少了参数数量。

然而，如果坐标系不是按照 DH 约定放置的，则变换矩阵可能无法直接表示为上述形式，需要先调整坐标系。
  
\end{sol}

\question 

\begin{sol}

设第一次重合时的变换矩阵为 $ _W^0\mathbf{T}_1 $，第二次重合时的变换矩阵为 $ _W^0\mathbf{T}_2 $。这些变换矩阵可以通过正运动学计算得到。

变换矩阵可以写为矩阵分块的形式，即

\[
  _W^0\mathbf{T}_1 = 
  \begin{bmatrix}
    R_1 & t_1 \\
    0 & 1
  \end{bmatrix}
\]

\[
  _W^0\mathbf{T}_2 = 
  \begin{bmatrix}
    R_2 & t_2 \\
    0 & 1
  \end{bmatrix}
\]

两次变换工具点与固定点重合时，满足：
\[
  _W^0\mathbf{T}_1 {^6}P = _W^0\mathbf{T}_2 {^6}P
\]

其中 $ {^6}P = [{^6}p_x, {^6}p_y, {^6}p_z, 1]^T $ 为齐次坐标形式。展开得，
\[
  R_1 {^6}P + t_1 = R_2 {^6}P + t_2
\]

整理得，
\[
  (R_1 - R_2) {^6}P = t_2 - t_1
\]

因此，工具偏移量 $ {^6}P $ 的表达式为：
\[
 {^6}P = (R_1 - R_2)^{-1} (t_2 - t_1)
\]
（我们可以通过调整机械臂的姿态来保证矩阵 $ R_1 - R_2 $ 一定可逆）
  
\end{sol}

\question 

\begin{sol}

由于关节3的轴线方向发生了变化，所以需要对坐标系做一点改动。

\img<q8>[0.5]{Question 8}{figs/q8.png}

\tbl<tb8>{DH-table for question 8} {
    i & $ a_{i-1} $ & $ \alpha_{i-1} $ & $ d_i $ & $ \theta_i $ \\
    1 & 0 & 0 & $ a_1 $ & $ \theta_1 $ \\
    2 & 0 & $ -\frac{\pi}{2} $ & $ d_3 $ & $ \theta_2 $ \\
    3 & 0 & $ -\frac{\pi}{2} $ & $ a_2 + d $ & 0 \\
    4 & $ a_3 $ & 0 & $ d_4 $ & $ \theta_4 $ \\
    5 & 0 & $ \frac{\pi}{2} $ & 0 & $ \theta_5 $ \\
    6 & 0 & $ -\frac{\pi}{2} $ & 0 & $ \theta_6 $
  }

各连杆变换矩阵如下：

\[
  _1^0\mathbf{T} =
  \begin{bmatrix}
    \cos\theta_1 & -\sin\theta_1 & 0 & 0 \\
    \sin\theta_1 & \cos\theta_1 & 0 & 0 \\
    0 & 0 & 1 & a_1 \\
    0 & 0 & 0 & 1
  \end{bmatrix}
\]

\[
  _2^1\mathbf{T} =
  \begin{bmatrix}
    \cos\theta_2 & -\sin\theta_2 & 0 & 0 \\
    0 & 0 & 1 & d_3 \\
    -\sin\theta_2 & -\cos\theta_2 & 0 & 0 \\
    0 & 0 & 0 & 1
  \end{bmatrix}
\]

\[
  _3^2\mathbf{T} =
  \begin{bmatrix}
    1 & 0 & 0 & 0 \\
    0 & 0 & 1 & a_2 + d \\
    0 & -1 & 0 & 0 \\
    0 & 0 & 0 & 1
  \end{bmatrix}
\]

\[
  _4^3\mathbf{T} =
  \begin{bmatrix}
    \cos\theta_4 & -\sin\theta_4 & 0 & a_3 \\
    \sin\theta_4 & \cos\theta_4 & 0 & 0 \\
    0 & 0 & 1 & d_4 \\
    0 & 0 & 0 & 1
  \end{bmatrix}
\]

\[
  _5^4\mathbf{T} =
  \begin{bmatrix}
    \cos\theta_5 & -\sin\theta_5 & 0 & 0 \\
    0 & 0 & -1 & 0 \\
    \sin\theta_5 & \cos\theta_5 & 0 & 0 \\
    0 & 0 & 0 & 1
  \end{bmatrix}
\]

\[
  _6^5\mathbf{T} =
  \begin{bmatrix}
    \cos\theta_6 & -\sin\theta_6 & 0 & 0 \\
    0 & 0 & 1 & 0 \\
    -\sin\theta_6 & -\cos\theta_6 & 0 & 0 \\
    0 & 0 & 0 & 1
  \end{bmatrix}
\]

正运动学方程为：

\[
  _6^0\mathbf{T} = _1^0\mathbf{T} _2^1\mathbf{T} _3^2\mathbf{T} _4^3\mathbf{T} _5^4\mathbf{T} _6^5\mathbf{T} =
  \begin{bmatrix} 
r_{11} & r_{12} & r_{13} & p_x \\ 
r_{21} & r_{22} & r_{23} & p_y \\ 
r_{31} & r_{32} & r_{33} & p_z \\ 
0 & 0 & 0 & 1
\end{bmatrix}
\] 

\begin{align*} 
r_{11} &= c_{6}\,\left(c_{5}\,\left(s_{1}\,s_{4}+c_{1}\,c_{2}\,c_{4}\right)-c_{1}\,s_{2}\,s_{5}\right)+s_{6}\,\left(c_{4}\,s_{1}-c_{1}\,c_{2}\,s_{4}\right) \\ 
r_{12} &= c_{6}\,\left(c_{4}\,s_{1}-c_{1}\,c_{2}\,s_{4}\right)-s_{6}\,\left(c_{5}\,\left(s_{1}\,s_{4}+c_{1}\,c_{2}\,c_{4}\right)-c_{1}\,s_{2}\,s_{5}\right) \\ 
r_{13} &= -s_{5}\,\left(s_{1}\,s_{4}+c_{1}\,c_{2}\,c_{4}\right)-c_{1}\,c_{5}\,s_{2} \\ 
p_x &= a_{3}\,c_{1}\,c_{2}-c_{1}\,s_{2}\,\left(a_{2}+d\right)-d_{3}\,s_{1}-d_{4}\,c_{1}\,s_{2} \\ 
r_{21} &= -c_{6}\,\left(c_{5}\,\left(c_{1}\,s_{4}-c_{2}\,c_{4}\,s_{1}\right)+s_{1}\,s_{2}\,s_{5}\right)-s_{6}\,\left(c_{1}\,c_{4}+c_{2}\,s_{1}\,s_{4}\right) \\ 
r_{22} &= s_{6}\,\left(c_{5}\,\left(c_{1}\,s_{4}-c_{2}\,c_{4}\,s_{1}\right)+s_{1}\,s_{2}\,s_{5}\right)-c_{6}\,\left(c_{1}\,c_{4}+c_{2}\,s_{1}\,s_{4}\right) \\ 
r_{23} &= s_{5}\,\left(c_{1}\,s_{4}-c_{2}\,c_{4}\,s_{1}\right)-c_{5}\,s_{1}\,s_{2} \\ 
p_y &= d_{3}\,c_{1}-s_{1}\,s_{2}\,\left(a_{2}+d\right)+a_{3}\,c_{2}\,s_{1}-d_{4}\,s_{1}\,s_{2} \\ 
r_{31} &= s_{2}\,s_{4}\,s_{6}-c_{6}\,\left(c_{2}\,s_{5}+c_{4}\,c_{5}\,s_{2}\right) \\ 
r_{32} &= s_{6}\,\left(c_{2}\,s_{5}+c_{4}\,c_{5}\,s_{2}\right)+c_{6}\,s_{2}\,s_{4} \\ 
r_{33} &= c_{4}\,s_{2}\,s_{5}-c_{2}\,c_{5} \\ 
p_z &= a_{1}-c_{2}\,\left(a_{2}+d\right)-d_{4}\,c_{2}-a_{3}\,s_{2}
\end{align*}

\end{sol}

\question 

\begin{sol}

人在行走时，如果只考虑腿在竖直平面内的运动，那么就可以等价为一个平面三连杆机械臂。

\img<q9>[0.5]{Question 9}{figs/q9.png}

\tbl<tb9>{DH-table for question 9} {
    i & $ a_{i-1} $ & $ \alpha_{i-1} $ & $ d_i $ & $ \theta_i $ \\
    1 & 0 & 0 & 0 & $ \theta_{hip} $ \\
    2 & $ l_{thigh} $ & 0 & 0 & $ \theta_{knee} $ \\
    3 & $ l_{shin} $ & 0 & 0 & $ \theta_{ankle} + \frac{\pi}{2} $ \\
  }

首先计算出正运动学变换矩阵：
\[
  _3^0\mathbf{T} = _1^0\mathbf{T} _2^1\mathbf{T} _3^2\mathbf{T} =
  \begin{bmatrix} 
    -s_{123} & -c_{123} & 0 & l_{shin} c_{12} + l_{thigh} c_1 \\ 
    c_{123} & -s_{123} & 0 & l_{shin} s_{12} + l_{thigh} s_1 \\ 
    0 & 0 & 1 & 0 \\ 
    0 & 0 & 0 & 1
  \end{bmatrix}
\] 
其中 $ \theta_{hip} $ 简写为 $ \theta_1 $，$ \theta_{knee} $ 简写为 $ \theta_2 $，$ \theta_{ankle} $ 简写为 $ \theta_3 $。

在脚趾离地和脚跟触地的过程中，我们假设0号坐标系不动。

脚趾触地时，根据各关节角度，计算得脚踝的变换矩阵为
\[
 _3^0\mathbf{T}_1 = 
 \begin{bmatrix} 
    0.7061 & -0.7081 & 0 & 794.3064 \\ 
    0.7081 & 0.7061 & 0 & -305.6470 \\ 
    0 & 0 & 1 & 0 \\ 
    0 & 0 & 0 & 1
  \end{bmatrix}
\]

脚趾在脚踝坐标系中的位置向量为 $ ^{3}p_{toe} = (150, 0, 0, 1)^T $，所以脚趾在0号坐标系中的位置向量
\[
 ^{0}p_{toe} = _3^0\mathbf{T}{_1} ^{3}p_{toe} = 
  \begin{pmatrix} 
    900.2242 \\
    -199.4330 \\
    0 \\
    1.0000 \\
  \end{pmatrix}
\]
 
脚跟离地时，脚踝的变换矩阵为
\[
 _3^0\mathbf{T}_2 = 
 \begin{bmatrix} 
    -0.2001 & -0.9798 & 0 & 886.5435 \\ 
    0.9798 & -0.2001 & 0 & 12.4824 \\ 
    0 & 0 & 1 & 0 \\ 
    0 & 0 & 0 & 1
  \end{bmatrix}
\]

脚跟在脚踝坐标系中的位置向量为 $ ^{3}p_{heel} = (-50, 0, 0, 1)^T $，所以脚跟在0号坐标系中的位置向量
\[
 ^{0}p_{heel} = _3^0\mathbf{T}{_2} ^{3}p_{heel} = 
 \begin{pmatrix} 
    896.5461 \\
    -36.5068 \\
    0 \\
    1.0000
  \end{pmatrix}
\]

步幅是脚趾触地点与脚跟离地点之间的距离，也就是 $ ^{0}p_{toe} $ 和 $ ^{0}p_{heel} $ 之间的距离。我们忽略x方向分量的变化，只取y方向分量的变化，也就是
\[
  stride = ^{0}_{y}p_{heel} - ^{0}_{y}p_{toe} = 162.9261
\]
  
\end{sol}

\question 

\begin{sol}

\img<q10>[0.75]{Question 10}{figs/q10.png}

\tbl<tb10>{DH-table for question 10} {
    i & $ a_{i-1} $ & $ \alpha_{i-1} $ & $ d_i $ & $ \theta_i $ \\
    1 & 0 & 0 & $ d_1 $ & $ \theta_1 $ \\
    2 & 0 & $ \frac{\pi}{2} $ & 0 & $ \theta_2 $ \\
    3 & $ \sqrt{d_2^2 + d_3^2} $ & 0 & $ d_4 + x $ & 0 \\
    4 & 0 & 0 & 0 & $ \theta_4 $ \\
    5 & 0 & $ \frac{\pi}{2} $ & 0 & $ \theta_5 $ \\
    6 & 0 & $ -\frac{\pi}{2} $ & 0 & $ \theta_6 $ \\
    7 & 0 & 0 & y & 0 \\
  }

求出各坐标系变换矩阵如下：

\[
  _1^0\mathbf{T} =
  \begin{bmatrix}
    \cos\theta_1 & -\sin\theta_1 & 0 & 0 \\
    \sin\theta_1 & \cos\theta_1 & 0 & 0 \\
    0 & 0 & 1 & d_1 \\
    0 & 0 & 0 & 1
  \end{bmatrix}
\]

\[
  _2^1\mathbf{T} =
  \begin{bmatrix}
    \cos\theta_2 & -\sin\theta_2 & 0 & 0 \\
    0 & 0 & -1 & 0 \\
    \sin\theta_2 & \cos\theta_2 & 0 & 0 \\
    0 & 0 & 0 & 1
  \end{bmatrix}
\]

\[
  _3^2\mathbf{T} =
  \begin{bmatrix}
    1 & 0 & 0 & \sqrt{d_2^2 + d_3^2} \\
    0 & 1 & 0 & 0 \\
    0 & 0 & 1 & d_4 + x \\
    0 & 0 & 0 & 1
  \end{bmatrix}
\]

\[
  _4^3\mathbf{T} =
  \begin{bmatrix}
    \cos\theta_4 & -\sin\theta_4 & 0 & 0 \\
    \sin\theta_4 & \cos\theta_4 & 0 & 0 \\
    0 & 0 & 1 & 0 \\
    0 & 0 & 0 & 1
  \end{bmatrix}
\]

\[
  _5^4\mathbf{T} =
  \begin{bmatrix}
    \cos\theta_5 & -\sin\theta_5 & 0 & 0 \\
    0 & 0 & -1 & 0 \\
    \sin\theta_5 & \cos\theta_5 & 0 & 0 \\
    0 & 0 & 0 & 1
  \end{bmatrix}
\]

\[
  _6^5\mathbf{T} =
  \begin{bmatrix}
    \cos\theta_6 & -\sin\theta_6 & 0 & 0 \\
    0 & 0 & 1 & 0 \\
    -\sin\theta_6 & -\cos\theta_6 & 0 & 0 \\
    0 & 0 & 0 & 1
  \end{bmatrix}
\]

\[
  _7^6\mathbf{T} =
  \begin{bmatrix}
    1 & 0 & 0 & 0 \\
    0 & 1 & 0 & 0 \\
    0 & 0 & 1 & y \\
    0 & 0 & 0 & 1
  \end{bmatrix}
\]

最终的正运动学变换矩阵为：

\[
  _7^0\mathbf{T} = _1^0\mathbf{T} _2^1\mathbf{T} _3^2\mathbf{T} _4^3\mathbf{T} _5^4\mathbf{T} _6^5\mathbf{T} _7^6\mathbf{T} =
  \begin{bmatrix} 
r_{11} & r_{12} & r_{13} & p_x \\ 
r_{21} & r_{22} & r_{23} & p_y \\ 
r_{31} & r_{32} & r_{33} & p_z \\ 
0 & 0 & 0 & 1
\end{bmatrix}
\] 

\begin{align*} 
r_{11} &= c_{6}\,\left(s_{1}\,s_{5}-c_{5}\,\left(c_{1}\,s_{2}\,s_{4}-c_{1}\,c_{2}\,c_{4}\right)\right)-s_{6}\,\left(c_{1}\,c_{2}\,s_{4}+c_{1}\,c_{4}\,s_{2}\right) \\ 
r_{12} &= -s_{6}\,\left(s_{1}\,s_{5}-c_{5}\,\left(c_{1}\,s_{2}\,s_{4}-c_{1}\,c_{2}\,c_{4}\right)\right)-c_{6}\,\left(c_{1}\,c_{2}\,s_{4}+c_{1}\,c_{4}\,s_{2}\right) \\ 
r_{13} &= c_{5}\,s_{1}+s_{5}\,\left(c_{1}\,s_{2}\,s_{4}-c_{1}\,c_{2}\,c_{4}\right) \\ 
p_x &= s_{1}\,x+d_{4}\,s_{1}+y\,\left(c_{5}\,s_{1}+s_{5}\,\left(c_{1}\,s_{2}\,s_{4}-c_{1}\,c_{2}\,c_{4}\right)\right)+c_{1}\,c_{2}\,\sqrt{{d_{2}}^2+{d_{3}}^2} \\ 
r_{21} &= -c_{6}\,\left(c_{1}\,s_{5}+c_{5}\,\left(s_{1}\,s_{2}\,s_{4}-c_{2}\,c_{4}\,s_{1}\right)\right)-s_{6}\,\left(c_{2}\,s_{1}\,s_{4}+c_{4}\,s_{1}\,s_{2}\right) \\ 
r_{22} &= s_{6}\,\left(c_{1}\,s_{5}+c_{5}\,\left(s_{1}\,s_{2}\,s_{4}-c_{2}\,c_{4}\,s_{1}\right)\right)-c_{6}\,\left(c_{2}\,s_{1}\,s_{4}+c_{4}\,s_{1}\,s_{2}\right) \\ 
r_{23} &= s_{5}\,\left(s_{1}\,s_{2}\,s_{4}-c_{2}\,c_{4}\,s_{1}\right)-c_{1}\,c_{5} \\ 
p_y &= \left(-c_{1}\right)\,x+c_{2}\,s_{1}\,\sqrt{{d_{2}}^2+{d_{3}}^2}-y\,\left(c_{1}\,c_{5}-s_{5}\,\left(s_{1}\,s_{2}\,s_{4}-c_{2}\,c_{4}\,s_{1}\right)\right)-d_{4}\,c_{1} \\ 
r_{31} &= c_{24}\,s_{6}+s_{24}\,c_{5}\,c_{6} \\ 
r_{32} &= c_{24}\,c_{6}-s_{24}\,c_{5}\,s_{6} \\ 
r_{33} &= -s_{24}\,s_{5} \\ 
p_z &= d_{1}+s_{2}\,\sqrt{{d_{2}}^2+{d_{3}}^2}-y\,s_{24}\,s_{5}
\end{align*}

当机械手处于平衡状态时，各关节变量的值如下：
\begin{align*} 
  \theta_1 &= \frac{\pi}{2} \\
  \theta_2 &= -\arctan(\frac{d_3}{d_2}) \\
  x &= 0 \\
  \theta_4 &= \frac{\pi}{2} + \arctan(\frac{d_3}{d_2}) \\
  \theta_5 &= 0 \\
  \theta_6 &= 0 \\
  y &= 0 \\
\end{align*}

代入计算得机械手平衡状态时的变换矩阵为
\[
\begin{bmatrix}
    0 & 0 & 1 & d_4 \\
    0 & -1 & 0 & d_2 \\
    1 & 0 & 0 & d_1 - d_3 \\
    0 & 0 & 0 & 1
\end{bmatrix}
\]

我们可以从另一个角度得到这个变换矩阵。直接看图可知，机械手末端相对于基坐标系的位移为 $ ^0\mathbf{p}_{7O} = (d_4, d_2, d_1 - d_3)^T $。再看机械手末端坐标系相对于基坐标系的旋转关系，可以写出旋转矩阵为
\[
^0_7\mathbf{R} = 
\begin{bmatrix}
  0 & 0 & 1 \\
  0 & -1 & 0 \\
  1 & 0 & 0
\end{bmatrix}
\]

因此，机械手平衡状态时的变换矩阵为
\[
\begin{bmatrix}
  ^0_7\mathbf{R} & ^0\mathbf{p}_{7O} \\
  0 & 1
\end{bmatrix}
\]

这印证了我们之前计算结果的正确。

\end{sol}

\question

\begin{sol}

图中原有坐标系不符合我们建系的习惯，遂重新建立坐标系如下：

\img<q11>[0.75]{Question 11}{figs/q11.png}

\tbl<tb11>{DH-table for question 11} {
    i & $ a_{i-1} $ & $ \alpha_{i-1} $ & $ d_i $ & $ \theta_i $ \\
    1 & 0 & 0 & $ l_0 $ & $ \theta_1 $ \\
    2 & 0 & $ -\frac{\pi}{2} $ & $ l_1 $ & $ \theta_2 $ \\
    3 & 0 & $ -\frac{\pi}{2} $ & $ l_2 $ & $ \theta_3 $ \\
    4 & $ l_3 $ & 0 & 0 & $ \theta_4 $ \\
    5 & $ l_4 $ & 0 & $ l_5 $ & $ \theta_5 $ \\
    6 & 0 & $ -\frac{\pi}{2} $ & $ l_6 $ & $ \theta_6 $ \\
    7 & 0 & $ -\frac{\pi}{2} $ & $ l_7 $ & $ \theta_7 $ \\
  }

求出各坐标系变换矩阵如下：

\[
  _1^0\mathbf{T} =
  \begin{bmatrix}
    \cos\theta_1 & -\sin\theta_1 & 0 & 0 \\
    \sin\theta_1 & \cos\theta_1 & 0 & 0 \\
    0 & 0 & 1 & l_0 \\
    0 & 0 & 0 & 1
  \end{bmatrix}
\]

\[
  _2^1\mathbf{T} =
  \begin{bmatrix}
    \cos\theta_2 & -\sin\theta_2 & 0 & 0 \\
    0 & 0 & 1 & l_1 \\
    -\sin\theta_2 & -\cos\theta_2 & 0 & 0 \\
    0 & 0 & 0 & 1
  \end{bmatrix}
\]

\[
  _3^2\mathbf{T} =
  \begin{bmatrix}
    \cos\theta_3 & -\sin\theta_3 & 0 & 0 \\
    0 & 0 & 1 & l_2 \\
    -\sin\theta_3 & -\cos\theta_3 & 0 & 0 \\
    0 & 0 & 0 & 1
  \end{bmatrix}
\]

\[
  _4^3\mathbf{T} =
  \begin{bmatrix}
    \cos\theta_4 & -\sin\theta_4 & 0 & l_3 \\
    \sin\theta_4 & \cos\theta_4 & 0 & 0 \\
    0 & 0 & 1 & 0 \\
    0 & 0 & 0 & 1
  \end{bmatrix}
\]

\[
  _5^4\mathbf{T} =
  \begin{bmatrix}
    \cos\theta_5 & -\sin\theta_5 & 0 & l_4 \\
    \sin\theta_5 & \cos\theta_5 & 0 & 0 \\
    0 & 0 & 1 & l_5 \\
    0 & 0 & 0 & 1
  \end{bmatrix}
\]

\[
  _6^5\mathbf{T} =
  \begin{bmatrix}
    \cos\theta_6 & -\sin\theta_6 & 0 & 0 \\
    0 & 0 & 1 & l_6 \\
    -\sin\theta_6 & -\cos\theta_6 & 0 & 0 \\
    0 & 0 & 0 & 1
  \end{bmatrix}
\]

\[
  _7^6\mathbf{T} =
  \begin{bmatrix}
    \cos\theta_7 & -\sin\theta_7 & 0 & 0 \\
    0 & 0 & 1 & l_7 \\
    -\sin\theta_7 & -\cos\theta_7 & 0 & 0 \\
    0 & 0 & 0 & 1
  \end{bmatrix}
\]

在计算 $ _7^8\mathbf{T} $ 时，我们只需要以坐标系{$ \mathbf{B}_4 $ }作为参考坐标系即可。
\[
  _7^4\mathbf{T} = _5^4\mathbf{T} _6^5\mathbf{T} _7^6\mathbf{T} =
  \begin{bmatrix} 
    s_{5}s_{7} + s_{5}s_{6}s_{7} & s_{7}s_{5} - s_{5}s_{6}s_{7} & -s_{5}s_{6} & l_4 - l_6s_{5} - l_7s_{5}s_{6} \\ 
    s_{6}s_{7}s_{5} - s_{5}s_{7} & -s_{5}s_{7} - s_{6}s_{5}s_{7} & -s_{5}s_{6} & l_6s_{5} - l_7s_{5}s_{6} \\ 
    -s_{7}s_{6} & s_{6}s_{7} & -s_{6} & l_5 - l_7s_{6} \\ 
    0 & 0 & 0 & 1
  \end{bmatrix}
\]

从 $ _7^4\mathbf{T} $ 中提取出坐标系{$ \mathbf{B}_7 $ }的原点在坐标系{$ \mathbf{B}_4 $ }中的表示 $ ^{4}p_{7O} $为：
\[
  ^{4}p_{7O} = 
  \begin{pmatrix}
    l_4 - l_6s_{5} - l_7s_{5}s_{6} \\
    l_6s_{5} - l_7s_{5}s_{6} \\
    l_5 - l_7s_{6}
  \end{pmatrix}
\]

根据题意，坐标系{ $ \mathbf{B}_8 $ }的原点在坐标系{ $ \mathbf{B}_4 $ }中的表示 $ ^{4}p_{8O} $可以写出：
\[
  ^{4}p_{8O} = 
  \begin{pmatrix}
    l_4 - a \\
    b \\
    0
  \end{pmatrix}
\]

因为 $ x_8 $ 轴指向坐标系{ $ \mathbf{B}_7 $ }的原点，所以
\[
  ^{4}X_{8} = ^{4}p_{7O} - ^{4}p_{8O} =
  \begin{pmatrix}
    a - l_6s_{5} - l_7s_{5}s_{6} \\
    l_6s_{5} - l_7s_{5}s_{6} - b \\
    l_5 - l_7s_{6}
  \end{pmatrix}
\]

将其单位化，得
\[
  ^{4}\hat{X}_{8} =
  \frac{1}{\sqrt{\left(a - l_6 s_5 - l_7 s_5 s_6\right)^2 + \left(l_6 s_5 - l_7 s_5 s_6 - b\right)^2 + \left(l_5 - l_7 s_6\right)^2}}
  \begin{pmatrix}
    a - l_6 s_5 - l_7 s_5 s_6 \\
    l_6 s_5 - l_7 s_5 s_6 - b \\
    l_5 - l_7 s_6
  \end{pmatrix}
\]

由题意知，$ ^{4}\hat{Z}_{8} $ 即为 $ (0, 0, 1)^T $。因为x轴与z轴必须正交，所以
\[
  ^{4}\hat{X}_{8} \cdot ^{4}\hat{Z}_{8} = l_5 - l_7s_{6} = 0
\]

即 $ s_{6} = \frac{l_5}{l_7} $。要想使解存在，则 $ l_5 $ 必须小于等于 $ l_7 $。$ \theta_6 $ 存在2个解，分别是 $ \arcsin{\frac{l_5}{l_7}} $ 和 $ \pi - \arcsin{\frac{l_5}{l_7}} $。

$ ^{4}\hat{Y}_{8} $ 可由 $ ^{4}\hat{Z}_{8} $ 和 $ ^{4}\hat{X}_{8} $ 叉乘得到：
\[
  ^{4}\hat{Y}_{8} = ^{4}\hat{Z}_{8} \times ^{4}\hat{X}_{8} = 
  \frac{1}{\sqrt{\left(a - l_6 s_5 - l_7 s_5 s_6\right)^2 + \left(l_6 s_5 - l_7 s_5 s_6 - b\right)^2 + \left(l_5 - l_7 s_6\right)^2}}
  \begin{pmatrix}
    l_6s_{5} - l_7s_{5}s_{6} - b \\
    l_6s_{5} + l_7s_{5}s_{6} - a \\
    0
  \end{pmatrix}
\]

所以我们得到了变换矩阵
\[
  _8^4\mathbf{T} = 
  \begin{bmatrix} 
    \frac{a - l_6 s_5 - l_5 s_5}{\sqrt{(a - l_6 s_5 - l_5 s_5)^2 + (l_6 s_5 - l_5 s_5 - b)^2}} & \frac{l_6 s_5 - l_5 s_5 - b}{\sqrt{(a - l_6 s_5 - l_5 s_5)^2 + (l_6 s_5 - l_5 s_5 - b)^2}} & 0 & l_4 - a \\
    \frac{l_6 s_5 - l_5 s_5 - b}{\sqrt{(a - l_6 s_5 - l_5 s_5)^2 + (l_6 s_5 - l_5 s_5 - b)^2}} & \frac{l_6 s_5 + l_5 s_5 - a}{\sqrt{(a - l_6 s_5 - l_5 s_5)^2 + (l_6 s_5 - l_5 s_5 - b)^2}} & 0 & b \\
    0 & 0 & 1 & 0 \\
    0 & 0 & 0 & 1
  \end{bmatrix}
\]

所以其他矩阵也就迎刃而解了：
\[
  _7^8\mathbf{T} = _8^4\mathbf{T}^{-1} {_7}^4\mathbf{T} = 
\]

\[
  \small
  \begin{bmatrix} 
    \frac{A\,\left(s_{5}\,s_{7}+c_{5}\,c_{6}\,c_{7}\right)}{\sqrt{A^2+B^2}}-\frac{B\,\left(c_{5}\,s_{7}-c_{6}\,c_{7}\,s_{5}\right)}{\sqrt{A^2+B^2}} & 
    \frac{A\,\left(c_{7}\,s_{5}-c_{5}\,c_{6}\,s_{7}\right)}{\sqrt{A^2+B^2}}-\frac{B\,\left(c_{5}\,c_{7}+c_{6}\,s_{5}\,s_{7}\right)}{\sqrt{A^2+B^2}} & 
    -\frac{s_{6}\,\left(A\,c_{5}+B\,s_{5}\right)}{\sqrt{A^2+B^2}} & 
    -\frac{B\,b-A\,a-B\,l_{6}\,c_{5}+A\,l_{6}\,s_{5}+A\,l_{7}\,c_{5}\,s_{6}+B\,l_{7}\,s_{5}\,s_{6}}{\sqrt{A^2+B^2}} \\[8pt] 
    \frac{A\,\left(c_{5}\,s_{7}-c_{6}\,c_{7}\,s_{5}\right)}{\sqrt{A^2+B^2}}+\frac{B\,\left(s_{5}\,s_{7}+c_{5}\,c_{6}\,c_{7}\right)}{\sqrt{A^2+B^2}} & 
    \frac{A\,\left(c_{5}\,c_{7}+c_{6}\,s_{5}\,s_{7}\right)}{\sqrt{A^2+B^2}}+\frac{B\,\left(c_{7}\,s_{5}-c_{5}\,c_{6}\,s_{7}\right)}{\sqrt{A^2+B^2}} & 
    -\frac{s_{6}\,\left(B\,c_{5}-A\,s_{5}\right)}{\sqrt{A^2+B^2}} & 
    \frac{A\,b+B\,a-A\,l_{6}\,c_{5}-B\,l_{6}\,s_{5}-B\,l_{7}\,c_{5}\,s_{6}+A\,l_{7}\,s_{5}\,s_{6}}{\sqrt{A^2+B^2}} \\[8pt] 
    -c_{7}\,s_{6} & s_{6}\,s_{7} & -c_{6} & l_{5}-l_{7}\,c_{6} \\[6pt] 
    0 & 0 & 0 & 1
  \end{bmatrix}
\]

为了表示方便，我们将 $ l_6 s_5 - l_5 s_5 - b $ 和 $ a - l_5 s_5 - l_6 s_5 $ 分别简写为 $ B $ 和 $ A $。

\[
  _8^3\mathbf{T} = _4^3\mathbf{T} _8^4\mathbf{T} =
  \begin{bmatrix} 
\frac{A\,c_{4}-B\,s_{4}}{\sqrt{A^2+B^2}} & \frac{B\,c_{4}+A\,s_{4}}{\sqrt{A^2+B^2}} & 0 & l_{3}-b\,s_{4}-c_{4}\,\left(a-l_{4}\right) \\ 
\frac{B\,c_{4}+A\,s_{4}}{\sqrt{A^2+B^2}} & -\frac{A\,c_{4}-B\,s_{4}}{\sqrt{A^2+B^2}} & 0 & b\,c_{4}-s_{4}\,\left(a-l_{4}\right) \\ 
0 & 0 & 1 & 0 \\ 
0 & 0 & 0 & 1
\end{bmatrix}
\]

在这里需要注意，我们可以看到关节3和关节4是平行的，所以 $ _3^2\mathbf{T} $ 和 $ _4^3\mathbf{T} $ 的乘积用和差化积公式将得到一个简化的表达式。
\[
  _4^2\mathbf{T} = _3^2\mathbf{T} _4^3\mathbf{T} =
  \begin{bmatrix} 
    c_{34} & -s_{34} & 0 & l_{3}\,c_{3} \\ 
    0 & 0 & 1 & l_{2} \\ 
    -s_{34} & -c_{34} & 0 & -l_{3}\,s_{3} \\ 
    0 & 0 & 0 & 1
  \end{bmatrix}
\]

\[
  _8^2\mathbf{T} = _4^2\mathbf{T} _8^4\mathbf{T} =
  \begin{bmatrix} 
\frac{A\,c_{34}-B\,s_{34}}{\sqrt{A^2+B^2}} & \frac{B\,c_{34}+A\,s_{34}}{\sqrt{A^2+B^2}} & 0 & l_{3}\,c_{3}-c_{34}\,\left(a-l_{4}\right)-b\,s_{34} \\ 
0 & 0 & 1 & l_{2} \\ 
-\frac{B\,c_{34}+A\,s_{34}}{\sqrt{A^2+B^2}} & \frac{A\,c_{34}-B\,s_{34}}{\sqrt{A^2+B^2}} & 0 & s_{34}\,\left(a-l_{4}\right)-b\,c_{34}-l_{3}\,s_{3} \\ 
0 & 0 & 0 & 1
\end{bmatrix}
\]

\[
  _8^1\mathbf{T} = _2^1\mathbf{T} _8^2\mathbf{T} =
  \begin{bmatrix} 
\frac{c_{2}\,\left(A\,c_{34}-B\,s_{34}\right)}{\sqrt{A^2+B^2}} & \frac{c_{2}\,\left(B\,c_{34}+A\,s_{34}\right)}{\sqrt{A^2+B^2}} & -s_{2} & -c_{2}\,\left(b\,s_{34}+c_{34}\,\left(a-l_{4}\right)-l_{3}\,c_{3}\right)-l_{2}\,s_{2} \\ 
-\frac{B\,c_{34}+A\,s_{34}}{\sqrt{A^2+B^2}} & \frac{A\,c_{34}-B\,s_{34}}{\sqrt{A^2+B^2}} & 0 & l_{1}-b\,c_{34}+s_{34}\,\left(a-l_{4}\right)-l_{3}\,s_{3} \\ 
-\frac{s_{2}\,\left(A\,c_{34}-B\,s_{34}\right)}{\sqrt{A^2+B^2}} & -\frac{s_{2}\,\left(B\,c_{34}+A\,s_{34}\right)}{\sqrt{A^2+B^2}} & -c_{2} & s_{2}\,\left(b\,s_{34}+c_{34}\,\left(a-l_{4}\right)-l_{3}\,c_{3}\right)-l_{2}\,c_{2} \\ 
0 & 0 & 0 & 1
\end{bmatrix}
\]

\[
  _8^0\mathbf{T} = _1^0\mathbf{T} _8^1\mathbf{T} =
  \begin{bmatrix} 
r_{11} & r_{12} & r_{13} & p_x \\ 
r_{21} & r_{22} & r_{23} & p_y \\ 
r_{31} & r_{32} & r_{33} & p_z \\ 
0 & 0 & 0 & 1
\end{bmatrix}
\] 

\begin{align*} 
r_{11} &= s_{1}\,\left(\frac{B\,c_{34}}{\sqrt{A^2+B^2}}+\frac{A\,s_{34}}{\sqrt{A^2+B^2}}\right)+c_{1}\,c_{2}\,\left(\frac{A\,c_{34}}{\sqrt{A^2+B^2}}-\frac{B\,s_{34}}{\sqrt{A^2+B^2}}\right) \\ 
r_{12} &= c_{1}\,c_{2}\,\left(\frac{B\,c_{34}}{\sqrt{A^2+B^2}}+\frac{A\,s_{34}}{\sqrt{A^2+B^2}}\right)-s_{1}\,\left(\frac{A\,c_{34}}{\sqrt{A^2+B^2}}-\frac{B\,s_{34}}{\sqrt{A^2+B^2}}\right) \\ 
r_{13} &= -c_{1}\,s_{2} \\ 
p_x &= -c_{1}\,\left(c_{2}\,\left(b\,s_{34}+c_{34}\,\left(a-l_{4}\right)-l_{3}\,c_{3}\right)+l_{2}\,s_{2}\right)-s_{1}\,\left(l_{1}-b\,c_{34}+s_{34}\,\left(a-l_{4}\right)-l_{3}\,s_{3}\right) \\ 
r_{21} &= c_{2}\,s_{1}\,\left(\frac{A\,c_{34}}{\sqrt{A^2+B^2}}-\frac{B\,s_{34}}{\sqrt{A^2+B^2}}\right)-c_{1}\,\left(\frac{B\,c_{34}}{\sqrt{A^2+B^2}}+\frac{A\,s_{34}}{\sqrt{A^2+B^2}}\right) \\ 
r_{22} &= c_{1}\,\left(\frac{A\,c_{34}}{\sqrt{A^2+B^2}}-\frac{B\,s_{34}}{\sqrt{A^2+B^2}}\right)+c_{2}\,s_{1}\,\left(\frac{B\,c_{34}}{\sqrt{A^2+B^2}}+\frac{A\,s_{34}}{\sqrt{A^2+B^2}}\right) \\ 
r_{23} &= -s_{1}\,s_{2} \\ 
p_y &= c_{1}\,\left(l_{1}-b\,c_{34}+s_{34}\,\left(a-l_{4}\right)-l_{3}\,s_{3}\right)-s_{1}\,\left(c_{2}\,\left(b\,s_{34}+c_{34}\,\left(a-l_{4}\right)-l_{3}\,c_{3}\right)+l_{2}\,s_{2}\right) \\ 
r_{31} &= -\frac{s_{2}\,\left(A\,c_{34}-B\,s_{34}\right)}{\sqrt{A^2+B^2}} \\ 
r_{32} &= -\frac{s_{2}\,\left(B\,c_{34}+A\,s_{34}\right)}{\sqrt{A^2+B^2}} \\ 
r_{33} &= -c_{2} \\ 
p_z &= l_{0}+s_{2}\,\left(b\,s_{34}+c_{34}\,\left(a-l_{4}\right)-l_{3}\,c_{3}\right)-l_{2}\,c_{2}
\end{align*}

\end{sol}

\end{document}
